%% ---------------------------------------------------------------------------
\section{Implementation cost against algorithm cost}
\label{sec:backends}

This section exists because the measurements forced it, and it went wrong once
before it went right.

\subsection{How the question arose}

An early comparison showed ML-DSA-65 signing on the M1 Pro about five times
slower than on a modest x86 server, while SHA-256, SHA-3 and ECDSA ran faster
on the M1 Pro, which ruled out throttling, the efficiency cores and Keccak.
Through liboqs on both machines most of the gap closed, and through the
Python package's OpenSSL the two libraries were 47 times apart on the M1 Pro.
A study that stopped there would blame OpenSSL's arm64 code path. That
attribution is wrong.

\subsection{Every build we could obtain}

\begin{table}[t]
\caption{ML-DSA-65 through every implementation and build we could obtain on the M1 Pro, and the corresponding builds on the Xeon VM, 1,000 timed calls each after 50 discarded. The C harness uses one call sequence for every OpenSSL build, the same one the Python \texttt{cryptography} package uses. The last two rows are the same package version over the same OpenSSL version and differ only in whether that OpenSSL was compiled with optimisation.}
\label{tab:impls}
\small
\setlength{\tabcolsep}{4.5pt}
\begin{tabular}{llrrrr}
\toprule
Implementation & Build & Sign, med. & Sign, mean & Verify & $\times$ fastest\\
\midrule
\multicolumn{6}{l}{\emph{Apple M1 Pro, arm64}}\\
liboqs 0.16.0 & native, arm64 path & 0.097 & 0.117 & 0.047 & 1.0\\
liboqs 0.16.0 & generic, no NEON & 0.181 & 0.226 & 0.065 & 1.9\\
PQClean reference & portable C, \texttt{-O3} & 0.314 & 0.398 & 0.102 & 3.2\\
OpenSSL 4.0.2 & source, \texttt{-O3} & 0.477 & 0.599 & 0.110 & 4.9\\
OpenSSL 4.0.2 & source, \texttt{-O3}, no asm & 0.563 & 0.757 & 0.133 & 5.8\\
OpenSSL 3.6.4 & Homebrew & 0.496 & 0.603 & 0.115 & 5.1\\
\texttt{cryptography} 50.0.1 & OpenSSL 4.0.2 source, \texttt{-O3} & 0.431 & 0.563 & 0.107 & 4.4\\
\texttt{cryptography} 50.0.1 & PyPI wheel, no \texttt{-O} & 4.599 & 6.416 & 0.859 & 47.2\\
\midrule
\multicolumn{6}{l}{\emph{Intel Xeon VM, x86-64, GCC 13.3}}\\
liboqs 0.16.0 & AVX2 & 0.137 & 0.171 & 0.039 & 1.0\\
liboqs 0.16.0 & generic & 0.358 & 0.479 & 0.138 & 2.6\\
PQClean reference & \texttt{-O3} & 0.573 & 0.748 & 0.316 & 4.2\\
OpenSSL 4.0.2 & source, asm & 0.801 & 1.036 & 0.182 & 5.8\\
OpenSSL 4.0.2 & source, no asm & 0.821 & 1.097 & 0.201 & 6.0\\
\bottomrule
\end{tabular}
\end{table}


Table~\ref{tab:impls} times ML-DSA-65 through every build we could obtain on
each machine, with one C harness driving every OpenSSL build through the call
sequence the Python package uses, so that between rows only the build
changes. Read the last two M1 rows first. They are the same
\texttt{cryptography} release over the same OpenSSL release, and one signs
ten times slower than the other. OpenSSL records the compiler line it was
built with, and the wheel's reads
\begin{quote}\small\ttfamily
cc -fPIC -arch arm64 -mmacosx-version-min=11.0 -DL\_ENDIAN\\
\quad -DOPENSSL\_PIC -D\_REENTRANT -DNDEBUG
\end{quote}
with no optimisation flag, so Clang compiles at \texttt{-O0}. A profile of the
wheel while it signs agrees: most samples land in
\texttt{constant\_time\_\allowbreak select\_32},
\texttt{reduce\_\allowbreak once} and \texttt{value\_\allowbreak barrier\_32},
helpers a few instructions long meant to be inlined into the number-theoretic
transform, and the wheel carries eight out-of-line copies of each. The Linux
x86-64 wheel of the same release records \texttt{-O3}, which accounts for the
difference that looked like an architecture effect.

The defect is easy to miss because it spares exactly the operations a sanity
check would use: SHA-2, SHA-3 and elliptic-curve arithmetic go through
hand-written assembly on arm64, which does not care about \texttt{-O}, and
every post-quantum algorithm in OpenSSL is C (Table~\ref{tab:buildeffect} in
the supplement). The same compiler line appears in every macOS wheel we
examined, from 3.4 in 2021 to 50.0.1. Its cause is \texttt{CFLAGS} set in the
build workflow's environment, which OpenSSL's \texttt{Configure} takes as a
replacement for its defaults; we reported it with a two-line fix, merged the
same day~\cite{pyca_issue15702,pyca_infra796}, and the macOS wheel of 50.0.2
records \texttt{-O3 -Wall}.

With the build held equal the picture becomes ordinary, and the two machines
agree. On the M1 Pro, correctly built implementations of ML-DSA-65 span a
factor of six: liboqs's NEON path~\cite{becker2022neon} is fastest, the same
library with vectorisation off costs 1.9 times as much, PQClean's portable
reference~\cite{kannwischer2022pqclean} 3.2 times, and OpenSSL 4.4 to 5.8
times. On the Xeon VM the spread is \NxSpread: liboqs's AVX2 path signs
\NavxSign{} times and verifies \NavxVerify{} times faster than its generic
build, and OpenSSL signs \NxOsslVsOqs{} times slower than AVX2 liboqs. Either
spread exceeds the spread across the three ML-DSA parameter sets through one
library, about 2.2, so implementation matters more than parameter set once the
build is sound. On the Pi~5's Cortex-A76, through a wheel built at
\texttt{-O3}, OpenSSL signs 5.7 times slower than liboqs, inside the same
range. Through the source-built OpenSSL with assembly, deterministic signing, which for a given key, message and context always takes the same rejection-sampling path, has a 95th percentile within six per cent of its median; randomised signing has one
three times its median, and Section~\ref{sec:resign} shows what that tail
does to a batch.

\subsection{What this asks of the field}
\label{sec:audit}

Both builds in the last two M1 rows report themselves as
\texttt{cryptography} 50.0.1 over OpenSSL 4.0.2, so a paper that named both
would still have published a figure ten times too slow and had no way to know.
A post-quantum timing needs the implementation, its version and the flags it
was compiled with. To see how often applied work reports them, we coded \NauditN{} studies published from 2022 to 2026 that time their own post-quantum
signatures or key encapsulation in an applied setting (TLS and QUIC, IoT and
embedded, VPN, DNSSEC, blockchain, vehicles, PKI), excluding papers whose
contribution is an optimised implementation, which always report their build
(Table~\ref{tab:audit}; the protocol is in Section~\ref{sec:auditprotocol} of the supplement). Of the \NauditN, \NauditLib{} name their library and \NauditVer{} its version; \NauditCc{} name a compiler and \NauditFlags{} say something about the build, but only \NauditOlevel{} state an explicit optimisation level such as \texttt{-O3}, and those same \NauditBoth{} are the only ones to give both compiler and optimisation level. Of the \NauditOqs{} that use liboqs, \NauditOqsBuild{} give the optimisation level they built it with and \NauditOqsAny{} say anything about the build, although its optimised and portable paths differ by a factor of two or three (Table~\ref{tab:impls}). \NauditNoSpread{} report no measure of spread, \NauditMeanNoSpread{} of them a mean alone and \NauditNoStat{} no statistic at all, so a reader cannot tell whether a difference exceeds the run-to-run spread. One author coded every study, without a second rater, and the sample leans to open-access and preprint work; the counts describe this sample, not the field.

The effect is not confined to signatures. Our matched-implementation control
on key encapsulation attributes a factor of 4.2 to 4.7 of the gap between HQC
and ML-KEM to implementation maturity~\cite{alfarajat2026hqc}. The control
costs little: run the candidate algorithms through one further
implementation, record how each was built, and report both, with a measure
of spread.
